\chapter{Zeitformen}
\label{cha:Zeitformen}

Die Zeitformen des deutschen Verbs ermöglichen die präzise zeitliche Einordnung von Handlungen.

\section{Präsens (Gegenwart)}
\label{sec:Praesens}

Das Präsens wird verwendet für:
\begin{itemize}
    \item \textbf{Gegenwart}: Ich arbeite heute.
    \item \textbf{Zukunft}: Morgen fliege ich nach Berlin.
    \item \textbf{historische Gegenwart}: 1789 bricht die Französische Revolution aus.
\end{itemize}

\section{Präteritum (Vergangenheit)}
\label{sec:Praeteritum}

Das Präteritum wird verwendet für:
\begin{itemize}
    \item \textbf{abgeschlossene Handlungen in der Vergangenheit}: Ich sah den Film gestern.
    \item \textbf{Erzählsprache}: Er ging nach Hause.
    \item \textbf{gewohnte Handlungen}: Als Kind ging er jeden Tag in die Schule.
\end{itemize}

\section{Perfekt (Vergangenheit)}
\label{sec:Perfekt}

Das Perfekt wird verwendet für:
\begin{itemize}
    \item \textbf{abgeschlossene Handlungen in der gesprochenen Sprache}: Ich habe den Film gesehen.
    \item \textbf{heute abgeschlossene Handlungen}: Heute habe ich viel gearbeitet.
\end{itemize}

\subsection{Perfekt mit Modalverben}
\label{sec:PerfektModal}

Bei Modalverben gibt es zwei Formen:

\begin{tcolorbox}[title=Standardform (Doppelter Infinitiv)]
\hfill\newline
haben/sein + Infinitiv (Modalverb) + Partizip II (Vollverb)
\end{tcolorbox}

\begin{itemize}
    \item Ich habe arbeiten müssen. (= Ich musste arbeiten.)
    \item Sie hat nicht kommen können. (= Sie konnte nicht kommen.)
    \item Er hat das Buch lesen wollen. (= Er wollte das Buch lesen.)
\end{itemize}

\begin{tcolorbox}[title=Ersatzform]
\hfill\newline
haben/sein + zu + Infinitiv + Partizip II (gehabt)
\end{tcolorbox}

\begin{itemize}
    \item Ich habe zu arbeiten gehabt. (= Ich musste arbeiten.)
    \item Wir haben früher zu gehen gehabt. (= Wir mussten früher gehen.)
\end{itemize}

\section{Plusquamperfekt (Vorvergangenheit)}
\label{sec:Plusquamperfekt}

Bildung: hatte/waren + Partizip II

\begin{itemize}
    \item Nachdem sie die Prüfung bestanden hatte, feierte sie.
    \item Als ich ankam, waren sie schon gegangen.
    \item Bevor er anrief, hatte ich ihn gesucht.
\end{itemize}

\section{Futur I (Zukunft)}
\label{sec:FuturI}

Bildung: werden + Infinitiv

\begin{itemize}
    \item \textbf{Zukunft}: Morgen werde ich arbeiten.
    \item \textbf{Vermutung}: Er wird etwa 40 Jahre alt sein.
    \item \textbf{Willen/Befehl}: Du wirst das tun!
\end{itemize}

\section{Futur II (Zukunft + Abgeschlossen)}
\label{sec:FuturII}

Bildung: werden + Partizip II + haben/sein

\begin{itemize}
    \item \textbf{abgeschlossene Zukunft}: Nächstes Jahr werde ich das Examen gemacht haben.
    \item \textbf{Vermutung über Vergangenheit}: Er wird schon gegangen sein.
\end{itemize}

\section{Übungen zu Zeitformen}
\label{sec:ZeitformenUebungen}

\subsection{Perfekt mit Modalverben}
\begin{enumerate}
    \item „Ich darf das nicht tun." $\rightarrow$ Perfekt \quad \textit{(Lösung: Ich habe das nicht tun dürfen.)}
    \item „Sie kann nicht kommen." $\rightarrow$ Perfekt \quad \textit{(Lösung: Sie hat nicht kommen können.)}
    \item „Er will das Buch lesen." $\rightarrow$ Perfekt \quad \textit{(Lösung: Er hat das Buch lesen wollen.)}
    \item „Wir müssen früher gehen." $\rightarrow$ Ersatzform \quad \textit{(Lösung: Wir haben früher zu gehen gehabt.)}
    \item „Ich soll die Hausaufgaben machen." $\rightarrow$ Perfekt \quad \textit{(Lösung: Ich habe die Hausaufgaben machen sollen.)}
    \item „Sie darf nicht rauchen." $\rightarrow$ Perfekt \quad \textit{(Lösung: Sie hat nicht rauchen dürfen.)}
    \item „Er kann gut schwimmen." $\rightarrow$ Perfekt \quad \textit{(Lösung: Er hat gut schwimmen können.)}
\end{enumerate}

\subsection{Plusquamperfekt}
\begin{enumerate}
    \item Nachdem sie die Prüfung \underline{\hspace{1cm}}, feierte sie. (bestehen) \quad \textit{(Lösung: hatte bestanden)}
    \item Als ich ankam, \underline{\hspace{1cm}} sie. (gehen) \quad \textit{(Lösung: war gegangen)}
    \item Bevor er anrief, \underline{\hspace{1cm}} ich ihn. (suchen) \quad \textit{(Lösung: hatte gesucht)}
    \item Sie war traurig, weil ihr Hund \underline{\hspace{1cm}}. (sterben) \quad \textit{(Lösung: war gestorben)}
    \item Ich \underline{\hspace{1cm}} den Zug, deshalb kam ich zu spät. (verpassen) \quad \textit{(Lösung: hatte verpasst)}
    \item Als du klingeltest, \underline{\hspace{1cm}} ich. (duschen) \quad \textit{(Lösung: hatte geduscht)}
    \item Wir \underline{\hspace{1cm}} noch nie in Berlin, bevor wir reisten. (sein) \quad \textit{(Lösung: waren gewesen)}
\end{enumerate}

\subsection{Futur}
\begin{enumerate}
    \item Morgen \underline{\hspace{1cm}} ich. (reisen) \quad \textit{(Lösung: werde reisen)}
    \item Das \underline{\hspace{1cm}} teuer. (sein) \quad \textit{(Lösung: wird teuer sein)}
    \item Nächstes Jahr \underline{\hspace{1cm}} ich das Examen. (machen) \quad \textit{(Lösung: werde das Examen gemacht haben)}
    \item Wenn du kommst, \underline{\hspace{1cm}} ich schon. (gehen) \quad \textit{(Lösung: werde ich schon gegangen sein)}
    \item \underline{\hspace{1cm}} du pünktlich? (sein) \quad \textit{(Lösung: Wirst du pünktlich sein?)}
    \item Um 20 Uhr \underline{\hspace{1cm}} wir essen. (fertig sein) \quad \textit{(Lösung: werden wir essen fertig sein)}
    \item \underline{\hspace{1cm}} sie das Buch schon? (lesen) \quad \textit{(Lösung: Wird sie das Buch schon gelesen haben?)}
\end{enumerate}

\chapterend